\documentclass[12pt,oneside,a4paper]{article}

\usepackage[utf8]{inputenc} % Lærer LaTeX at forstå unicode - HUSK at filen skal
% være unicode (UTF-8), standard i Linux, ikke i
% Win.

%\usepackage[danish]{babel} % Så der fx står Figur og ikke Figure, Resumé og ikke
% Abstract etc. (god at have).

%\usepackage{graphicx}
\usepackage{amsfonts}
\usepackage{amsthm}        % Theorems
\usepackage{amsmath}
%\usepackage{hyperref}

%\renewcommand{\mid}[1]{{\rm E}\!\left[#1\right]}
\newcommand{\bas}{\begin{eqnarray*}}
\newcommand{\eas}{\end{eqnarray*}}
\newcommand{\be}{\begin{equation}}
\newcommand{\ee}{\end{equation}}
\newcommand{\bea}{\begin{eqnarray}}
\newcommand{\eea}{\end{eqnarray}}

%\newtheorem{thm}{Sætning}[section]
%\newtheorem{mydef}[thm]{Definition}
%\newtheorem{eks}[thm]{Eksempel}

\title{Solving quintics}
\date{December 2025}
\author{Michael Jørgensen}

\begin{document}

\maketitle

Character table for $S_5$:
{https://www2.math.upenn.edu/~chai/notes_clc_courses/course_notes_ugrad/charTableS_5-A_5.pdf}

Conjugacy classes:
{https://en.wikipedia.org/wiki/Conjugacy_class}

Character table:
{https://en.wikipedia.org/wiki/Character_table}

Character table for $F_{20}$:
{http://sporadic.stanford.edu/Math122/lecture14.pdf}

Frobenius group of order 20:
{http://manu.amiot.free.fr/pdf/Articles/chords.pdf}

Solving solvable quintics:
{https://www.ams.org/journals/mcom/1991-57-195/S0025-5718-1991-1079014-X/S0025-5718-1991-1079014-X.pdf}

\end{document}


